% ---------------------------------------------------------------------------
% Section IV-B Results and Discussion and IV-C Analysis. Fifth draft,
% 2026-09-15, after the lab meeting: three research questions (success,
% speed, generalization) tied to what, when and how much; fewer numbers,
% the tables and figures carry the values; "trick" for the family and
% "setting" for its variants. Every count is verified against
% paired_results_all.csv and the per-episode records (final_verify3.py).
% Style: no em-dashes, no semicolons, no prose colons, no hedges, no
% footnotes.
% ---------------------------------------------------------------------------

\subsection{Results and Discussion}

We read Tables~\ref{tab:simplerenv-results} and~\ref{tab:libero-results}
and Figs.~\ref{fig:teaser} and~\ref{fig:tradeoff-widowx}
to~\ref{fig:consistency} through three questions, one for each part of the
control loop that a trick touches.
\begin{list}{}{\setlength{\labelwidth}{2.9em}\setlength{\labelsep}{0.4em}%
  \setlength{\leftmargin}{3.3em}\setlength{\itemsep}{0pt}%
  \setlength{\parsep}{0pt}\setlength{\topsep}{2pt}}
\item[\textbf{RQ-1:}] Does changing \emph{what} the policy sees or
\emph{how much} it computes ever raise success, and which tricks leave
success intact?
\item[\textbf{RQ-2:}] Which tricks, by changing \emph{when} the policy is
called or \emph{how much} it computes, make it faster, and what does each
unit of speed cost?
\item[\textbf{RQ-3:}] Does a trick that helps one policy help another?
\end{list}

\textbf{RQ-1.} In Fig.~\ref{fig:teaser} no bar rises more than 10 points
above the dotted line of the original, while the bars that fall 10 points or more
below it belong to action repeat, depth pruning and, on OpenVLA and
UniVLA, foveation. Tables~\ref{tab:simplerenv-results}
and~\ref{tab:libero-results} say the same for every environment, with
few and small green success cells against many and large red ones. Of the 110 best-setting cells, 22
lower success significantly and four raise it, and the four gains, on
CogACT WidowX, OpenVLA Fractal and UniVLA Goal, are of the size and
number that chance allows over this many comparisons, while most of the
drops exceed 10 points. Over all 286 settings six gains pass against
about seven expected by chance, and two of the four table gains sit on
CogACT WidowX, the one pair whose every depth and fusion setting rises,
which the seeds of Section~V should test. Guarded reuse and temporal
fusion never lower success significantly in any best cell, and guarded
reuse in none of its 66 settings, although
eleven of its twelve OpenVLA LIBERO settings lose a few points and
pooling those four suites would make the moderate setting a significant
loss. The other
three each lower it on several backbones, foveation on three, action
repeat and depth pruning on four each at their mildest setting.
\textbf{[A-1] No trick raises success. Guarded reuse and temporal fusion
leave it intact on every backbone, and foveation, action repeat and depth
pruning each cost success on about half of them.}

\textbf{RQ-2.} In Fig.~\ref{fig:tradeoff-widowx} the shaded corner,
faster and no worse, is empty on SpatialVLA and UniVLA and holds at most
four settings elsewhere, none more than a few points above the original
except depth pruning on CogACT. Action repeat is the only trick far to
the right, at 1.4 to 3.6 times the original's speed per step, and it is
also the trick that falls furthest, with significant drops on 8 of 22
backbone and environment pairs at two steps and 20 at four
(Fig.~\ref{fig:ablation-depth}). Depth pruning moves a little to the
right and, on some backbones, far down. One removed layer is faster than
the original beyond the latency floor on 13 of the 22 pairs and costs
success significantly on 8, and four layers are faster on 21 and cost
success on 14. Guarded reuse saves time only where its gates open, a
tenth at most on OpenVLA LIBERO and about 5 percent or less elsewhere.
Foveation makes six backbones slower by 2 to 17 percent and moves UniVLA,
whose blur is amortized over a chunk, by under 4 percent, and task-aware
fusion adds 7 to 17 percent to every call on OpenVLA, MiniVLA and UniVLA.
% 2026-09-15: depth1 speedup beyond the 1.6 percent floor of IV-A on 13 of 22 pairs (14 at the earlier 1.4 percent floor; CronusVLA Fractal at 1.5 percent moves out). Foveation per-step slowdown over all settings: 2.0 (OpenVLA Goal keep20) to 16.8 percent (CronusVLA WidowX keep50) on the six backbones other than UniVLA; UniVLA -0.7 to +3.2 percent, eight of ten settings within the floor.
\textbf{[A-2] Action repeat
is the fastest trick and the most harmful, depth pruning buys a small
speedup on most backbones at a success cost that the backbone sets, and
the tricks that touch what the policy sees buy no speed at all.}

\textbf{RQ-3.} In Fig.~\ref{fig:tradeoff-widowx} depth pruning sits above
the original on CogACT and far below it on MiniVLA and SpatialVLA, on the same
tasks and episodes. Foveation at the same keep ratio raises UniVLA and
lowers SmolVLA on the same LIBERO suite. Action repeat at two steps
lowers success on all six WidowX backbones, significantly on four, and
leaves the four Fractal
backbones where they were (Fig.~\ref{fig:ablation-depth}).
Over the 13 combinations of Fig.~\ref{fig:consistency} no setting is
above the original on more than six,
and the settings that most often leave it unchanged are the ones whose
gates rarely fire. \textbf{[A-3] A trick does not transfer.
Its sign is a joint property of the rule, the backbone and the
environment, so a gain on one combination says nothing about the next
without a matched test there.}

\subsection{Analysis}

\textbf{Episode length follows success.} Every failed episode ran to its
step cap, so Avg.\ Steps in the tables tracks the Success column and is
not an independent measure of efficiency.

\textbf{What a repeated call saves.} Per-step latency under action repeat
lands above the halving that the call count predicts, because the
simulator step is paid on every environment step and only the model
computation scales with the calls.

\textbf{What the per-episode record adds.} On UniVLA WidowX,
motion-entropy and task-aware fusion both lose 11 of 200 episodes net,
and only the first loss is significant, because the two settings flip
different episodes. A setting can report the original policy under its
own name. The UniVLA guarded reuse table cells fired on at most five
steps and match the original on every episode but one, and
conservative-adaptive fusion reused no patch on most steps of SpatialVLA
and UniVLA. We therefore report how often each gate fired and recommend
the same for any gated trick. Cells on CogACT, CronusVLA, MiniVLA and
OpenVLA WidowX carry the run-to-run noise of Section~IV-A, and
CronusVLA's three fusion cells are one run whose settings were not
recorded.
